\documentclass[10pt]{article}
\usepackage[T1]{fontenc}
\usepackage[utf8]{inputenc}
\usepackage[french]{babel}
\usepackage{amsmath,amssymb,mathrsfs,mathpazo}
\usepackage{geometry}
\geometry{a4paper,margin=2cm}
\usepackage{graphicx,array,tabularx,tabularray,colortbl,multirow,multicol,booktabs,enumitem,pifont,ulem,fancybox,diagbox,cancel}
\usepackage{pstricks,pstricks-add,pst-plot,pst-tree,pst-node,pst-text,pst-func,pst-math,pst-bspline,pst-3dplot,pst-eucl,pst-solides3d}
\let\pscalculate\undefined
\usepackage{pst-bezier}
\usepackage{tikz,pgfplots,tkz-tab}
\usetikzlibrary{arrows,arrows.meta,patterns,calc,positioning,trees,decorations.markings,shapes.geometric,shapes.misc,shapes.arrows,intersections,decorations.pathmorphing,decorations.pathreplacing,backgrounds,shadows,babel}
\pgfplotsset{compat=1.18}
\usepgfplotslibrary{fillbetween}
\usepackage[np]{numprint}
\usepackage{listings,makecell,pas-tableur,twoopt}
\let\barre\undefined
\usepackage[pstricks]{tablvar}
\definecolor{darkgreen}{RGB}{0,100,0}
\providecommand{\red}{\color{red}}
\providecommand{\blue}{\color{blue}}
\providecommand{\green}{\color{green}}
\providecommand{\black}{\color{black}}
\providecommand{\white}{\color{white}}
\providecommand{\gray}{\color{gray}}
\providecommand{\cyan}{\color{cyan}}
\providecommand{\magenta}{\color{magenta}}
\providecommand{\yellow}{\color{yellow}}
\providecommand{\violet}{\color{violet}}
\providecommand{\orange}{\color{orange}}
\providecommand{\euro}{\texteuro}
\providecommand{\EUR}{\texteuro}
\providecommand{\ieme}{\textsuperscript{e}}
\providecommand{\degres}{\ensuremath{^\circ}}
\providecommand{\degree}{\ensuremath{^\circ}}
\providecommand{\vv}[1]{\overrightarrow{#1}}
\providecommand{\Coord}[1]{#1}
\providecommand{\decosix}{}
\providecommand{\starredbullet}{\ensuremath{\bullet}}
\pagestyle{empty}
\setlength{\parindent}{0pt}
\usepackage{adjustbox}\DeclareUnicodeCharacter{2236}{:}\DeclareUnicodeCharacter{2190}{\ensuremath{\leftarrow}}
\providecommand{\R}{\mathbb{R}}
\providecommand{\N}{\mathbb{N}}
\providecommand{\Z}{\mathbb{Z}}
\providecommand{\Q}{\mathbb{Q}}
\providecommand{\D}{\mathbb{D}}
\providecommand{\C}{\mathbb{C}}
\providecommand{\e}{{\mathrm{e}}}
\providecommand{\vect}[1]{\overrightarrow{#1}}
\providecommand{\vectt}[1]{\overrightarrow{\mathrm{#1}}}
\providecommand{\vvt}[1]{\overrightarrow{\mathrm{#1}}}
\providecommand{\np}[2][]{{#2~#1}}
\providecommand{\vv}[1]{{\overrightarrow{#1}}}
\renewcommand{\d}{{\mathrm{d}}}
\providecommand{\degres}{{^\circ}}
\providecommand{\degre}{{^\circ}}
\providecommand{\blue}{\color{blue}}
\providecommand{\red}{\color{red}}
\providecommand{\enskip}{\quad}
\providecommand{\llbracket}{\lbrack\!\lbrack}
\providecommand{\rrbracket}{\rbrack\!\rbrack}
\providecommand{\euro}{{\text{€}}}
\providecommand{\no}{{\text{n°}}}
\providecommand{\interval}[2]{{[#1;#2]}}
\providecommand{\intervFO}[2]{{[#1;#2[}}
\providecommand{\dotfill}{\ldots}

\providecommand{\ds}{\displaystyle}
\providecommand{\textding}[1]{#1}
\providecommand{\Oij}{\(\mathrm{O};\vec{\imath},\vec{\jmath}\)}
\providecommand{\Oijk}{\(\mathrm{O};\vec{\imath},\vec{\jmath},\vec{k}\)}
\providecommand{\Ouv}{\(\mathrm{O};\vec{u},\vec{v}\)}
\providecommand{\intervFF}[2]{\ensuremath{[#1;#2]}}
\renewcommand{\euro}{\ifmmode\text{\texteuro}\else\texteuro\fi}
\DeclareUnicodeCharacter{00B0}{\ifmmode^{\circ}\else\textdegree\fi}
\DeclareUnicodeCharacter{20AC}{\ifmmode\text{\texteuro}\else\texteuro\fi}
% Énoncés : Éducation nationale. Transcription et corrigés : APMEP, Denis Vergès et contributeurs. Découpage et classement : CoopMaths, Sébastien Lozano, Éric Elter et contributeurs. Adaptation de la présentation : Maths Explorer.
% CC BY-SA : https://coopmaths.fr/www/about/
% https://coopmaths.fr/alea/static/eam/2026/tex/eam_specialite_2026_04_polynesie_automatismes_cor.tex
\begin{document}
\noindent\textit{APMEP / CoopMaths. CC BY-SA. Présentation adaptée par Maths Explorer.}\par\medskip
\subsection*{Question 1 \quad Réponse D}

Un produit de facteurs est nul si et seulement si l'un au moins de ses facteurs est nul :
\[(2x - 6)(x + 4) = 0 \iff 2x - 6 = 0 \ \text{ ou } \ x + 4 = 0 \iff x = 3 \ \text{ ou } \ x = -4.\]
D'où $S = \{-4~;~3\}$.

\subsection*{Question 2 \quad Réponse B}

Le dé est équilibré : on est en situation d'équiprobabilité. Les issues strictement
supérieures à 4 sont 5 et 6, soit 2 issues favorables sur 6 :
\[P = \dfrac{2}{6} = \dfrac{1}{3}.\]

\subsection*{Question 3 \quad Réponse C}

On applique $\dfrac{a^m}{a^n} = a^{m-n}$ avec $3 = 3^1$ :
\[\dfrac{3}{3^5} = 3^{1-5} = 3^{-4}.\]

\subsection*{Question 4 \quad Réponse D}

La droite $\mathrm{D}_4$ passe par les points de coordonnées $(-1 ; 0)$ et $(0 ; 1)$. Son ordonnée à
l'origine est $b = 1$ et son coefficient directeur vaut
$a = \dfrac{1 - 0}{0 - (-1)} = 1$. Son équation réduite est donc $y = x + 1$.

\subsection*{Question 5 \quad Réponse B}

Un point appartient à $\mathcal{C}_f$ lorsque son ordonnée est l'image de son abscisse.
Pour H$(-3~;~28)$ :
\[f(-3) = 2\times(-3)^2 - 3\times(-3) + 1 = 18 + 9 + 1 = 28.\]
Le point H appartient donc à $\mathcal{C}_f$.

\subsection*{Question 6 \quad Réponse B}

On réduit d'abord la parenthèse au même dénominateur :
\[E = \left(\dfrac{1}{2} + \dfrac{1}{3}\right)\times\dfrac{2}{5}
= \left(\dfrac{3}{6} + \dfrac{2}{6}\right)\times\dfrac{2}{5}
= \dfrac{5}{6}\times\dfrac{2}{5} = \dfrac{10}{30} = \dfrac{1}{3}.\]

\subsection*{Question 7 \quad Réponse C}

Une hausse de 10\,\% correspond au coefficient multiplicateur $1,10$ et une baisse de
20\,\% au coefficient $0,80$. Le coefficient multiplicateur global est le produit
des deux :
\[1,10 \times 0,80 = 0,88 = 1 - 0,12.\]
Le prix a subi une baisse de 12\,\%.

\subsection*{Question 8 \quad Réponse C}

Toutes les grandeurs sont strictement positives :
\[g = \dfrac{Gm}{r^2} \iff r^2 g = Gm \iff r^2 = \dfrac{Gm}{g} \iff r = \sqrt{\dfrac{Gm}{g}}.\]

\end{document}
